\chapter{Bounded Completeness: Is the Top Region Actually Necessary? (Day 10)}
\label{chap:day10}

\section{Motivation}

Completeness is one of the six metrics defined for evaluating black-box explainable-AI frameworks
\citep{arreche2024exai}: does the explanation a model offers actually account for its decision, or
does the decision survive even when the explanation is taken away? A full completeness analysis would
ablate the entire feature space and ask whether \emph{any} subset of removed evidence changes the
output --- a combinatorially large undertaking well outside this pilot's scope. Day 10 instead
implements a deliberately \textbf{bounded} variant, scoped down to the one region the pipeline already
treats as most important. Per \texttt{technical-limitations-take-and-mitigation.md}:

\begin{quote}
\itshape ``A sample passes the bounded completeness check if the model provides a usable explanation
and if perturbing the top-ranked region or top two regions causes a measurable change in the model's
answer, confidence proxy, or grounding output.''
\end{quote}

This is a necessity test, not a sufficiency test: it asks whether the single region
Chapter~\ref{chap:day04}'s \texttt{standardize\_regions} ranked first is actually load-bearing for the
model's answer, or whether the model would have reached the same conclusion regardless of whether that
region existed at all. Chapters~\ref{chap:day04} and~\ref{chap:day05} already built every piece this
chapter needs --- ranked regions and the masked-rerun answer-change signal --- so Day 10's task is
substantially a rollup of work already done, plus one narrowly bounded extra rerun set to check a risk
the rollup alone cannot resolve.

\section{Implementation}

\texttt{completeness.py}'s \texttt{classify\_sample()} is a pure rollup over two existing columns: Day
4's \texttt{region\_extraction.csv} (\texttt{top\_region\_source}, which records whether a real
model-returned box or only the grid fallback was available) and Day 5's
\texttt{descriptive\_accuracy.csv} (\texttt{answer\_changed\_top1} and \texttt{answer\_changed\_top2}).
No new model calls are required for the headline classification: a sample with no model-returned region
at all is \texttt{no\_usable\_explanation}; a sample with a model region whose masking changed the
answer (top-1 or top-2) is \texttt{explanation\_supported}; a sample with a model region whose masking
did \emph{not} change the answer is \texttt{explanation\_weak}. \texttt{scripts/10\_bounded\_completeness.py}
ran this rollup over all 163 pilot samples end-to-end and wrote \texttt{results/bounded\_completeness.csv}
(163 rows).

One extra, explicitly bounded set of fresh model calls was made on top of the rollup: \textbf{159
top-1-only mask reruns}, reusing Chapter~\ref{chap:day04}'s already-extracted top-1 box rather than
re-ranking from scratch, that additionally capture post-mask \texttt{worker\_boxes}. This was necessary
because Chapter~5's own CSV never logged post-mask boxes, only the answer-change outcome; the rollup
alone cannot tell whether an answer change reflects the top region genuinely mattering or reflects the
same worker-detection-fallback-rerouting artifact that Chapters~\ref{chap:day05} and~\ref{chap:day09}
already documented being triggered by other means. Resolving that question --- the substance of
Finding~2 below --- required this one additional, narrowly scoped rerun set, not a new pipeline stage.

\textbf{Why no confidence-proxy threshold.} The bounded definition explicitly permits a measurable
change in the confidence proxy as an alternative passing signal alongside answer change. This was
checked before being ruled out: among Chapter~5's \texttt{answer\_changed\_top1 == False} rows, the
distribution of \texttt{abs(confidence\_drop\_top1)} (median 0.032, 75th percentile 0.057) is
essentially indistinguishable from the same statistic over the full population (median 0.032, 75th
percentile 0.057). Confidence drop carries no discriminative signal between "answer changed" and
"answer did not change" at this scale; adding a threshold on top of it would inject noise into the
classification rather than recover real signal. \texttt{classify\_sample()} therefore uses only the
categorical answer-change signal, consistent with Chapters~\ref{chap:day03} and~\ref{chap:day09}'s
standing caution that this pilot's confidence figure is an uncalibrated proxy, not a correctness signal.

\section{Results}

Table~\ref{tab:day10-headline} reports the headline classification across all 163 samples.

\begin{center}
\begin{tabular}{p{4.6cm}p{1.6cm}p{1.6cm}}
\toprule
\textbf{Verdict} & \textbf{n} & \textbf{\%} \\
\midrule
\texttt{explanation\_weak} & 88 & 54.0\% \\
\texttt{explanation\_supported} & 71 & 43.6\% \\
\texttt{no\_usable\_explanation} & 4 & 2.5\% \\
\bottomrule
\end{tabular}
\captionof{table}{Headline bounded-completeness classification, all 163 pilot samples.}
\label{tab:day10-headline}
\end{center}

A slim majority of samples (54.0\%) land in \texttt{explanation\_weak}: the top-ranked region exists and
is masked, but the answer does not move. Read uncritically, that could be mistaken for "the model is
mostly ignoring its own top-ranked evidence" --- Findings~1 and~2 below show that this headline number
is itself shaped by a fixable mechanism, not an inherent property of the model.

\textbf{Hard subsets are reported separately here, by deliberate policy, never blended into the overall
rate}, exactly as the plan specifically asks. Table~\ref{tab:day10-hardsubsets} reports two such subsets
alongside the overall figure for comparison.

\begin{center}
\begin{tabular}{p{4.0cm}p{2.4cm}p{2.0cm}p{1.6cm}p{1.9cm}}
\toprule
\textbf{Subset} & \textbf{n} & \textbf{supported} & \textbf{weak} & \textbf{no\_\allowbreak usable} \\
\midrule
overall & 163 & 43.6\% & 54.0\% & 2.5\% \\
\texttt{quality\_of\_info == "poor info"} & 84 (51.5\% of pilot) & 39.3\% & 58.3\% & 2.4\% \\
multi-rule-violation (n=7, too small for conclusions) & 7 & 42.9\% & 57.1\% & 0.0\% \\
\bottomrule
\end{tabular}
\captionof{table}{Hard subsets reported separately from the overall rate, by deliberate policy, never
blended into the headline figure.}
\label{tab:day10-hardsubsets}
\end{center}

The poor-info subset --- just over half the pilot --- is mildly worse than the overall rate (39.3\% vs.
43.6\% supported) but not dramatically so; whatever is driving \texttt{explanation\_weak}'s headline
share is not concentrated in the samples with poor underlying information quality. The
multi-rule-violation subset has only 7 samples, too few to support any real conclusion; it happens to
land close to the overall rate, but with $n=7$ that proximity is not itself evidence of anything.

\section{Finding 1: The Area-Ranking Mechanism, Again}

The headline 54.0\% \texttt{explanation\_weak} rate is not uniform across rules, and the pattern it
follows is the same mechanism Chapter~\ref{chap:day04} flagged as a foreseen limitation back when
\texttt{standardize\_regions} was first designed: ranking candidate regions by box area, descending,
rather than by relevance to the rule being tested. Table~\ref{tab:day10-byrule} breaks the verdict down
by rule.

\begin{center}
\begin{tabular}{p{5.0cm}p{2.0cm}p{1.6cm}p{1.9cm}p{1.1cm}}
\toprule
\textbf{Rule} & \textbf{supported} & \textbf{weak} & \textbf{no\_\allowbreak usable} & \textbf{n} \\
\midrule
rule\_1 (hard hat / vest) & \textbf{33.3\%} & 66.7\% & 0.0\% & 63 \\
rule\_2 (harness / lanyard) & 46.2\% & 50.0\% & 3.8\% & 26 \\
rule\_3 (guardrail / edge barrier) & \textbf{55.1\%} & 38.8\% & 6.1\% & 49 \\
rule\_4 (excavator proximity) & 44.0\% & 56.0\% & 0.0\% & 25 \\
\bottomrule
\end{tabular}
\captionof{table}{Bounded-completeness verdict by rule. rule\_1 is the weakest; rule\_3 is the
strongest --- a spread of nearly 22 percentage points in \texttt{explanation\_supported}.}
\label{tab:day10-byrule}
\end{center}

rule\_1 (hard hat/vest, the PPE rule the pilot cares about most directly) supports the model's answer
in only 33.3\% of cases; rule\_3 (guardrail/edge barrier) supports it in 55.1\% --- a 21.8 percentage
point spread between the weakest and strongest rule. This is not noise: it traces to exactly which
entity \texttt{standardize\_regions} ranks top-1 for each rule. For rule\_1, the worker's whole-body box
is almost always larger than the hard-hat or vest box it is being compared against, so the area-ranking
heuristic promotes the \textbf{worker} to top-1 in \textbf{48/50 (96\%)} of \texttt{ppe\_violation}/
rule\_1 samples. For rule\_3, the guardrail itself is frequently a large structural object spanning much
of the frame, so the \textbf{guardrail} is ranked top-1 in \textbf{33/50 (66\%)} of
\texttt{fall\_hazard}/rule\_3 samples, with the worker ranked top-1 in only 12/50 (24\%) of those same
samples.

Collapsing this across all 163 samples by which entity type was ranked top-1, rather than by rule, makes
the mechanism explicit, in Table~\ref{tab:day10-top1type}.

\begin{center}
\begin{tabular}{p{4.2cm}p{2.0cm}p{1.6cm}p{1.2cm}}
\toprule
\textbf{Top-1 region is\ldots} & \textbf{supported} & \textbf{weak} & \textbf{n} \\
\midrule
the worker & \textbf{37.5\%} & 62.5\% & 80 \\
the safety/hazard object & \textbf{49.4\%} & 45.8\% & 83 \\
\bottomrule
\end{tabular}
\captionof{table}{\texttt{explanation\_supported} rate by which entity type was ranked top-1. (Row
percentages for the safety/hazard-object row sum to 95.2\%; the remainder is
\texttt{no\_usable\_explanation}.)}
\label{tab:day10-top1type}
\end{center}

Masking the worker is reliably less likely to flip the answer than masking the actual safety or hazard
object --- 37.5\% versus 49.4\%, an 11.9 percentage point gap that holds across the whole pilot, not
just within one rule. \textbf{This is a real, fixable limitation of the ranking heuristic, not a
property of the model being explained.} Area is a reasonable proxy for visual salience in the absence
of any true attention or saliency signal from Florence-2 (the same constraint Chapter~\ref{chap:day04}
flagged at the point the heuristic was designed), but it systematically promotes the wrong entity to
top-1 exactly when the safety-relevant cue is small relative to the person wearing or standing near it
--- which is precisely the PPE case (rule\_1) this pilot, and the eventual paper, cares about most. A
ranking heuristic that instead ranked by "does this region match the rule's queried object class"
before falling back to area would directly correct the asymmetry visible in
Table~\ref{tab:day10-top1type}.

\section{Finding 2: A Worker-Fallback Artifact Inflates the Headline Number}

The 159 top-1-only mask reruns described in the Implementation section were run specifically to check
whether the worker-detection-fallback-rerouting risk --- flagged in Chapter~\ref{chap:day05} and
quantified under perturbation in Chapter~\ref{chap:day09} --- also contaminates this chapter's own
\texttt{explanation\_supported} count. It does, in a small but precisely quantifiable slice.

Of the 159 model-region samples, masking the top-1 region caused the rerun to lose its only detected
worker entirely (no worker redetected post-mask) in \textbf{12 cases (7.5\%)}. Of the 71 samples
classified \texttt{explanation\_supported}, \textbf{3 (4.2\%)} co-occur with this artifact: their
"supported" verdict is plausibly driven by \texttt{\_answer\_presence\_rule}'s scene-level fallback
branch (\texttt{"compliant" if object\_boxes else "violation"}) firing because no worker survived the
mask, rather than by the top region being genuinely necessary to the rule's actual reasoning. A
corrected, "genuinely-supported" count would be \textbf{68/163 (41.7\%)} rather than the naive
\textbf{71/163 (43.6\%)} reported as the headline figure above --- a small but real overstatement that
would have gone unflagged without this extra rerun set.

\textbf{Traced concretely on image \texttt{0000642} (rule\_2, baseline ``compliant'').}
Figure~\ref{fig:day10-finding2} shows the before/after pair. The model's "worker" box
\texttt{(436.2, 232.2, 499.8, 419.6)} and its "harness" box \texttt{(438.6, 232.2, 495.0, 419.6)} are
almost pixel-identical --- Florence-2 returned nearly the same region for both queries on this image.
Masking the top-ranked "worker" region therefore also erases the harness evidence in the same stroke,
visible in the after panel: the black box swallows the only worker-relevant region in frame. The rerun
fails to redetect any worker, triggers the fallback branch, and the answer changes --- registering as
\texttt{explanation\_supported}, when what actually happened is the same degenerate-fallback mechanism
first flagged in Chapter~\ref{chap:day05} and quantified under image perturbation in
Chapter~\ref{chap:day09}, here triggered by this chapter's own masking procedure rather than by noise
or a deliberate perturbation.

\begin{figure}[h]
\centering
\begin{subcaptionblock}{0.48\linewidth}
\centering
\includegraphics[width=\linewidth]{completeness/verify/0000642_before.png}
\caption{Before: the "worker" box and the "harness" box are almost pixel-identical --- Florence-2
returned nearly the same region for both queries.}
\end{subcaptionblock}
\hfill
\begin{subcaptionblock}{0.48\linewidth}
\centering
\includegraphics[width=\linewidth]{completeness/verify/0000642_after_mask.png}
\caption{After masking top-1 ("worker"): the harness evidence is erased in the same stroke; the rerun
finds no worker, the fallback branch fires, and the answer changes.}
\end{subcaptionblock}
\caption{Image \texttt{0000642}, rule~2: the worker-fallback artifact, traced. The "supported" verdict
this case earns reflects a degenerate fallback branch firing on a lost detection, not the top region
genuinely being necessary to the rule's reasoning.}
\label{fig:day10-finding2}
\end{figure}

Read in an aggregate table alone, \texttt{0000642} is indistinguishable from a genuine necessity result
--- masking the top region changed the answer, exactly what \texttt{explanation\_supported} is supposed
to mean. \textbf{It is not.} It is the same fallback mechanism that has now surfaced under three
different triggers across this pilot: deliberate masking (Chapter~\ref{chap:day05}), realistic
perturbation (Chapter~\ref{chap:day09}), and this chapter's own bounded-completeness mask reruns. Any
completeness-style metric built on top of a model with conditional or fallback logic needs to check for
this kind of artifact explicitly, the same way this chapter did with its extra 159-rerun set, rather
than trusting a raw answer-change rate to mean what it appears to mean.

\section{Finding 3: Grid-Fallback Is Rare but Concentrated in rule\_3}

Only \textbf{4/163 samples (2.5\%)} had no native grounding box at all and were classified
\texttt{no\_\allowbreak usable\_\allowbreak explanation}. Table~\ref{tab:day10-gridfallback} lists all
four.

\begin{center}
\begin{tabular}{p{2.0cm}p{2.6cm}p{1.6cm}p{2.6cm}}
\toprule
\textbf{Image ID} & \textbf{Primary class} & \textbf{Rule} & \textbf{Quality of info} \\
\midrule
\texttt{0000253} & fall\_hazard & rule\_3 & rich info \\
\texttt{0000361} & fall\_hazard & rule\_3 & poor info \\
\texttt{0000620} & compliant & rule\_2 & rich info \\
\texttt{0001229} & fall\_hazard & rule\_3 & poor info \\
\bottomrule
\end{tabular}
\captionof{table}{All four samples classified \texttt{no\_usable\_explanation} (no native grounding box
returned at all).}
\label{tab:day10-gridfallback}
\end{center}

Three of the four are rule\_3. "Guardrail" or "edge protection barrier" is a harder open-vocabulary
detection target for Florence-2 than "hard hat" or "worker," consistent with rule\_3 already showing the
lowest grounding hit-rate signal back in Chapters~\ref{chap:day03} and~\ref{chap:day04}. The rate is
small enough that it barely moves the headline number, but the concentration in a single rule is
consistent rather than random, and is worth noting alongside Finding~1's rule\_3 result: rule\_3 sits at
both extremes in this chapter, the strongest rule by \texttt{explanation\_supported} rate when a
grounding box \emph{is} returned, and the rule most likely to return no grounding box at all.

\section{Standard vs. Non-Standard Choices}

\textbf{Standard}: completeness, or necessity testing via top-$k$ region ablation, is itself standard
practice, directly the same metric defined for black-box XAI evaluation \citep{arreche2024exai}. The
``bounded'' scoping used here --- restricting the ablation to the top-ranked region or top two regions,
rather than the full feature space --- is an explicit, acknowledged narrowing of that standard
definition, adopted because a full completeness sweep over every combination of regions was out of
scope for a pilot of this size, not a departure from the metric's underlying logic.

\textbf{Non-standard}: two choices in this chapter are this pilot's own contribution rather than
inherited from the six-metric framework. First, the explicit policy of reporting hard subsets (poor-info
samples, multi-rule-violation samples) \textbf{separately, never blended into the overall rate} ---
Table~\ref{tab:day10-hardsubsets} rather than a single pooled number --- guards against a harder subset's
worse performance being diluted into an artificially reassuring headline figure, or conversely against a
too-small subset ($n=7$) being mistaken for a real signal. Second, and more significant: quantifying
\emph{exactly how much} the fallback-rerouting artifact inflates the headline number (Finding~2's
68/163 vs.\ 71/163 correction) via a dedicated extra rerun set, rather than reporting the naive rollup
number and trusting it. This is a rigor practice worth recommending generally for any completeness-style
metric built on top of a model that has conditional or fallback logic anywhere in its answer pipeline:
the metric's own perturbation procedure can itself trigger the same fallback paths that contaminate the
underlying model's answers, and a rollup alone cannot detect that without an extra, targeted check like
the one performed here.

\section{Implications for the Paper}

\textbf{Findings 1 and 2, taken together, make the strongest single case in this pilot for recommending
that the region-ranking heuristic be fixed before this pipeline is scaled up to a larger dataset.}
Finding~1 shows that the heuristic introduced in Chapter~\ref{chap:day04} --- rank candidate regions by
box area, descending --- produces a structural, rule-dependent bias: it promotes the worker to top-1 for
PPE rules (96\% of rule\_1 samples) precisely because the worker's box is almost always larger than the
PPE item itself, and masking the worker is measurably less likely to flip the answer (37.5\% supported)
than masking the actual safety object (49.4\% supported). That gap is not a property of Florence-2's
underlying competence; it is an artifact of which region the heuristic happens to rank first, and it
depresses the bounded-completeness rate hardest exactly on rule\_1, the PPE rule this pilot and the
eventual paper care about most directly. Finding~2 shows that the same underlying fallback mechanism
this pilot has now traced under three different triggers (Chapter~\ref{chap:day05}'s deliberate masking,
Chapter~\ref{chap:day09}'s realistic perturbations, and this chapter's own bounded mask reruns) also
inflates the headline \texttt{explanation\_supported} count by a small but precisely quantified amount
(3/71, correcting 43.6\% down to 41.7\%). Read together, the two findings say the same thing from two
directions: \textbf{the area-ranking heuristic is the single highest-leverage, most fixable problem
surfaced anywhere in this pilot}, because it is not a limitation of the underlying model or the
evaluation metric's design, but of one specific, identified, and replaceable function
(\texttt{standardize\_regions}) that every later masking-based chapter in this pilot --- Chapters~5, 6,
7, and this one --- inherits without modification.

This recommendation is not new to this chapter; Chapter~\ref{chap:day04} forward-referenced it at the
point the heuristic was first designed, anticipating that the area-based choice would resurface as a
limitation downstream. \textbf{This chapter is where that anticipated limitation is finally measured
directly and attributed to a specific, fixable cause}, and it should be presented as the paper's top
recommended next engineering step, tied directly to Chapter~14's task list item 3: replace
\texttt{standardize\_regions}'s area-only ranking with a rule-aware ranking that prefers a region
matching the rule's queried object class (hard hat, vest, harness, lanyard, guardrail, excavator) before
falling back to area only when no class match is available. Such a fix would not require any new model
calls or architectural change --- only a re-ranking of boxes Florence-2 has already returned --- and
Finding~1's Table~\ref{tab:day10-top1type} gives a direct, falsifiable prediction for what a successful
fix should produce: the 11.9 percentage point gap between masking the worker and masking the safety
object should narrow substantially, if not close outright, once the ranking stops systematically
promoting the wrong entity to top-1 for PPE rules.

Two secondary points should also carry forward. \textbf{First}, Finding~3's rule\_3 concentration in
\texttt{no\_usable\_explanation} is a second, independent piece of evidence (alongside
Chapters~\ref{chap:day03} and~\ref{chap:day04}'s already-noted grounding hit-rate signal) that
"guardrail"/"edge protection barrier" is a harder open-vocabulary target for Florence-2 than the other
rules' queried objects, worth flagging if rule\_3 is scaled up without also revisiting its grounding
prompt. \textbf{Second}, the Standard vs.\ Non-Standard section's rigor practice --- explicitly checking
whether a metric's own perturbation procedure can trigger the same fallback paths that contaminate the
underlying model, rather than assuming a clean separation between "the metric" and "the model under
test" --- generalizes beyond completeness and is worth stating in the paper's methods section as a
general caution for evaluating any model with conditional or fallback logic in its decision pipeline.
